\section{Introduction}
\label{sec:intro}

Multiple instance learning (MIL) is a weakly supervised setting in
which training data are grouped into \emph{bags} of instances and
labels are observed only at the bag level
\citep{dietterich1997solving,maron1998framework,andrews2003support}.
In the standard binary formulation, a bag is positive when it
contains at least one positive instance and negative otherwise; the
instance-level labels are latent. MIL is the natural model for
problems where annotation is cheap in aggregate but expensive per
instance: drug-activity prediction (a molecule is active if at least
one of its conformations binds), image classification (an image is
positive if some region contains the object of interest), and,
increasingly, computational pathology, where a gigapixel
whole-slide image carries a diagnosis but the diagnostic regions are
not delineated \citep{campanella2019clinical,lu2021data}.

A large methods literature attacks MIL: axis-parallel rectangles
\citep{dietterich1997solving}, Diverse Density
\citep{maron1998framework}, mi-SVM and MI-SVM
\citep{andrews2003support}, and attention-based deep MIL
\citep{ilse2018attention}, with problem variants surveyed by
\citet{carbonneau2018multiple} and \citet{foulds2010review}. These
methods differ in how they parametrize the instance-to-bag relation
and how they pool instance evidence. The gap between bag-level and
instance-level success is itself well documented:
\citet{vanwinckelen2016instance} show empirically that strong
bag-level accuracy need not transfer to instance-level accuracy, and
\citet{jang2024learnable} give a PAC account of when deep-MIL
algorithms are instance-level learnable. What these accounts do not
supply, and what we add, is a likelihood-identifiability reading in
the reliability (masked-cause / series-system) idiom: a closed-form
rank condition on the bag composition matrix that says \emph{when}
instance-type positivity is identifiable from bag labels alone, an OR-link
bag-prevalence identity that makes the bag-versus-instance gap of
\citet{vanwinckelen2016instance} a moment-matching statement, what
auxiliary information restores identifiability when the rank condition
fails, and what bias is incurred when the bag-formation assumption is
wrong. Practical questions follow: when are attention weights
interpretable as instance scores? When does a MIL model that predicts
bag labels well still mis-rank instances? How many instance-level
labels are needed to calibrate a bag-trained model? The closest prior
likelihood model is the EM treatment of MIL as Bernoulli MLE with
latent instance labels under the OR rule
\citep{chen2017milr}, which fits a logistic instance model by
expectation-maximization but does not give identifiability conditions
or a rank diagnostic; the closest identifiability-flavored neighbour
is partial-label learning \citep{cour2011learning}, in which each
\emph{instance} carries a candidate \emph{label} set with one
correct label, with MIL the dual coarsening in which a candidate
\emph{instance} set (the bag) carries one observed label. Both sit
inside the general coarse-data maximum-likelihood theory of
\citet{couso2017coarse}, of which the masked-cause specialization used
here is one case.

\subsection{The bridge to masked-data inference}

We observe that MIL is mathematically isomorphic to the masked-data
series-system identifiability problem of \citet{towell2026masked}.
In that framework, $m$ components in series produce a system that
fails when any one component fails; the failed component is masked,
but a \emph{candidate set} known to contain it is observed.
Identifiability conditions on the candidate-set matrix determine
when component-level parameters can be recovered.

The translation to MIL is direct. Instances play the role of
components. The bag plays the role of the candidate set. The ``at
least one positive instance'' rule is the series (logical-OR)
aggregation: a bag is positive exactly when the disjunction of its
instance labels is true, just as a series system fails exactly when
the disjunction of its component-failure events is true. A
\emph{singleton bag} (a bag with one instance, equivalently an
instance-level label) is a singleton candidate set. The
C1--C2--C3 coarsening conditions for ignorable masking port
directly. The framework was previously applied to scRNA-seq zero
inflation \citep{towell2026scrnacoarsening} and to
spatial-transcriptomics deconvolution
\citep{towell2026spatialcoarsening}; MIL is the discrete-label
member of the same family, with one structural inversion worth
flagging at the outset (\cref{sec:translation}): the \emph{negative}
bag is the fully informative observation (every instance label is
known to be $0$), while the \emph{positive} bag is the coarsened
one.

\subsection{Contributions}

We contribute:

\begin{itemize}
\item \textbf{A bridge from MIL to masked-data inference}
  (\cref{sec:translation}), formalizing mi-SVM, Diverse Density,
  attention-based MIL, and CLAM as special cases of one
  coarsening framework.

\item \textbf{A rank condition} (\cref{sec:identifiability}):
  instance-type positivity is identifiable from bag labels if and
  only if the bag-by-instance-type composition matrix has full
  column rank. Singleton bags contribute unit-vector rows and
  restore rank, the discrete-label analogue of how
  single-cell-resolution platforms restore identifiability in
  spatial transcriptomics.

\item \textbf{A bag-prevalence consistency theorem}
  (\cref{sec:identifiability}): any maximizer of the bag-label
  likelihood reproduces the observed bag-positive frequencies
  exactly along every direction in the column space of the
  composition matrix. Bag-level accuracy therefore cannot validate
  instance-level scores, the MIL counterpart of cell-total
  consistency in \citet{towell2026scrnacoarsening}.

\item \textbf{A bias bound under aggregation-rule misspecification}
  (\cref{sec:methodology}): a closed-form first-order bound on the
  instance-score bias incurred when the true bag rule is noisy-OR,
  label-noisy, or threshold rather than deterministic OR. The
  standard-versus-collective assumption taxonomy of
  \citet{foulds2010review} is recovered as a partition of
  aggregation rules by whether they preserve the C1 condition.

\item \textbf{A simulation that confirms the three theorems}
  (\cref{sec:validation}) with known instance-level ground truth,
  including a numerical demonstration of the firing-rate confound
  and of the contrast between weighted and unweighted bag-residual
  moment matching. A real-data application to MUSK1 and MUSK2
  (\cref{sec:musk}) demonstrates the rank diagnostic and the
  boundary-MLE regime on real chemical-descriptor data and quantifies
  the gap to published continuous-feature MIL baselines.
\end{itemize}

The central message: MIL methods are not arbitrary heuristics; they
are special cases of a single identifiability framework. The rank
condition is a diagnostic a practitioner can run before trusting an
instance-level output, attention weights in particular, since the
framework makes precise when those weights estimate a parameter and
when they merely interpolate a rank-deficient direction.
